\documentclass[10pt,a4paper]{amsart}
\usepackage[margin=19mm]{geometry}
\usepackage{amsmath,amssymb,mathtools}
\usepackage[scr=rsfs,cal=euler]{mathalfa}
\usepackage{xfrac}
\usepackage[unicode,hidelinks,bookmarks=false]{hyperref}
\newcommand{\ie}{i.e.\ }
\newcommand{\cf}{cf.\ }
\newcommand{\resp}{respectively\ }
\newcommand{\Subsection}{\subsection}
\providecommand{\nobreakdash}{\nobreak\hbox{-}}
\setlength{\emergencystretch}{2em}
\multlinegap0pt
\usepackage{bm}
\usepackage{bbm}
\usepackage{dsfont}
\usepackage{xspace}
\usepackage{cancel}
\usepackage{mathtools}
\usepackage{amsthm}
\makeatletter
\@ifundefined{theorem}{\newtheorem{theorem}{Theorem}}{}
\@ifundefined{defn}{\newtheorem{defn}[theorem]{Definition}}{}
\@ifundefined{lemma}{\newtheorem{lemma}[theorem]{Lemma}}{}
\makeatother
\newcommand\nph{\varphi}
\newcommand{\N}{\mathbb{N}}
\newcommand{\bb}[1]{\mathbb{#1}}
\newcommand{\cl}[1]{\mathcal{#1}}
\providecommand{\wt}[1]{\widetilde{#1}}
\title[An operator system approach to self-testing]{An operator system approach to self-testing}
\author{Jason Crann, Ivan G. Todorov, Lyudmila Turowska}
\date{Source: arXiv:2506.17980v1 (CC BY 4.0)}
\begin{document}
\maketitle

\noindent\emph{Source and licence.} An operator system approach to self-testing. Version arXiv:2506.17980v1 (CC BY 4.0), distributed under the Creative Commons Attribution 4.0 International licence, as recorded in the arXiv licence metadata for this exact version and shown on the abstract page https://arxiv.org/abs/2506.17980v1. Full rights notice: licenses/arxiv-2506.17980-CC-BY-4.0.txt.\par\smallskip

\noindent\textit{Editorial scope.} Counted unit: the Lemma whose source label is \texttt{p\_apextends} in the article (source0.tex lines 1376--1385, source label \texttt{p\_apextends}), delivered together with the article's own complete original proof of it (source0.tex lines 1387--1543, one \texttt{proof} environment of 157 lines). No mathematics is written, repaired or re-derived: statement and proof are the article's own text line for line. Required context retained uncounted: d\_appdil (lines 992--1026); eq\_supre (lines 1282--1284); uptoeps (lines 1299--1304); sep\_cycl (lines 1316--1324). Every definition, display and label of those lines is retained unchanged. Omitted from this package and not redistributed: 1--991, 1027--1281, 1285--1298, 1305--1315, 1325--1375, 1386--1386, 1544--4859. Nothing inside a retained excerpt is omitted or altered: every retained body is delivered exactly as the article's lines are written, so no omission receipt of any kind accompanies this package. Citations: the retained excerpts carry no citation command at all, so no bibliography entry of the article is transcribed and no reference is redistributed. Reference adaptation: none was needed and none was made; every reference inside the delivered unit resolves inside it, so no unresolved reference or citation is produced. Printed slips kept literal: no printed word, symbol, hypothesis or number of the counted statement or of its proof was altered. Layout and class: the article's own class is \texttt{amsart} with packages including latexsym, amssymb, amsmath, color, bbm, hyperref, lmodern, graphicx, tcolorbox, tabularx, tikz-cd; the wrapper is the lane's pinned \texttt{amsart} editorial wrapper at 10pt on A4, which restates the theorem-like environments the retained text needs and adds the article's own macro definitions that the retained text needs (\texttt{\textbackslash N}, \texttt{\textbackslash bb}, \texttt{\textbackslash cl}, \texttt{\textbackslash nph}, \texttt{\textbackslash wt}), with extra wrapper packages (bm, bbm, dsfont, xspace, cancel, mathtools). The delivered numbering is the unit's own. Assets: no retained span contains a figure environment, a graphics-inclusion command, a PDF-page inclusion, a table or a TikZ picture; no image file is copied into this package, the package declares no asset dependency and no image file is redistributed with it. Licence: version arXiv:2506.17980v1 of this article is distributed under the Creative Commons Attribution 4.0 International licence, as recorded in the arXiv licence metadata for this exact version and shown on the abstract page https://arxiv.org/abs/2506.17980v1. A full rights notice is delivered at licenses/arxiv-2506.17980-CC-BY-4.0.txt. Characters: every retained excerpt is pure ASCII, so the delivered engine, XeLaTeX with its default Latin Modern text font, reports no missing character.
\begin{defn}\label{d_appdil}
Let $\cl S_A$ and $\cl S_B$ be operator systems, and 
suppose that 
$S = (_{\cl A}H_{\cl B}, \varphi_A, \varphi_B, \xi)$  and 
$\wt{S} = (_{\wt{\cl A}}\wt{H}_{\wt{\cl B}}, \wt{\varphi}_A, \wt{\varphi}_B, \wt{\xi})$ are quantum commuting models over the pair $(\cl S_A,\cl S_B)$.

\begin{itemize}
\item[(i)] We say that $\wt{S}$
is a \emph{local dilation} of $S$, and write $S\preceq \wt{S}$, 
if there exist a bipartite quantum system 
$_{{\cl A}_{\rm aux}}(H_{\rm aux})_{{\cl B}_{\rm aux}}$ 
over a pair of von Neumann algebras (${\cl A}_{\rm aux}$, ${\cl B}_{\rm aux}$), a unit vector $\xi_{\rm aux}\in H_{\rm aux}$, and
a local isometry $V : \mbox{}_{\cl A}H_{\cl B}\to
\mbox{}_{\wt{\cl A}}\wt{H}_{\wt{\cl B}}\otimes
\mbox{}_{{\cl A}_{\rm aux}}(H_{\rm aux})_{{\cl B}_{\rm aux}}$ such that 
\begin{equation}\label{eq_TAB}
V\varphi_A(s) \varphi_B(t)\xi = (\tilde\varphi_A(s)\tilde\varphi_B(t)\tilde\xi) \otimes\xi_{\rm aux}, \ \ \ s\in \cl S_A, t\in\cl S_B.
\end{equation}

\item[(ii)] 
We say that $\wt{S}$
is an \emph{approximate local dilation} of $S$, and write $S\preceq_{\rm a} \wt{S}$, 
if there exist bipartite quantum systems
$_{{\cl A}_{i,\rm aux}}(H_{i,\rm aux})_{{\cl B}_{i,\rm aux}}$ 
over a pair of von Neumann algebras (${\cl A}_{i,\rm aux}$, ${\cl B}_{i,\rm aux}$), 
unit vectors $\xi_{i,\rm aux}\in H_{i,\rm aux}$, and
local isometries $V_i : \mbox{}_{\cl A}H_{\cl B}\to
\mbox{}_{\wt{\cl A}}\wt{H}_{\wt{\cl B}}\otimes
\mbox{}_{{\cl A}_{i,\rm aux}}(H_{i,\rm aux})_{{\cl B}_{i,\rm aux}}$, 
$i\in \bb{N}$, such that 
\begin{equation}\label{eq_TABa}
\left\|V_i\varphi_A(s) \varphi_B(t)\xi - (\tilde\varphi_A(s)\tilde\varphi_B(t)\tilde\xi) \otimes\xi_{i,\rm aux}\right\|\to_{i\to \infty} 0, \ \ \ s\in \cl S_A, t\in\cl S_B.
\end{equation}
\end{itemize}
\end{defn}

\begin{equation}\label{eq_supre}
p\xi = \xi.
\end{equation} 

\begin{lemma}\label{uptoeps}
    Let $\cl A\subseteq\cl B(H)$ be a von Neumann algebra and 
    $\xi\in H$ be a unit vector. 
    Assume that $\xi$ is separating for $\cl A$.
    Then, given $x\in\cl A$ and $\epsilon > 0$, there exists $y\in\cl A'$ such that $\|x\xi-y\xi\|<\epsilon$. 
\end{lemma}

\begin{lemma}\label{sep_cycl}
 Let $\cl A \subseteq \cl B(H)$ be a von Neumann algebra and  $\xi\in H$
 be a unit vector. Let  $\epsilon_A\in \cl A$ be the support 
 of the state $\omega_\xi|_{\cl A}$ of $\cl A$, 
 and $p_A\in \cl B(H)$ be the projection onto 
 $\overline{\cl A\xi}$.  Then $\xi \in \epsilon_Ap_AH$, and 
 $\xi$ is cyclic and separating for the von Neumann subalgebra 
 $\epsilon_Ap_A\cl A\epsilon_Ap_A$ of $\cl B(\epsilon_Ap_AH)$. 
\end{lemma}

\begin{lemma}\label{p_apextends}
 Let $S = (_{\cl A}H_{\cl B}, \varphi_A, \varphi_B, \xi)$ and
 $\wt{S} = (_{\wt{\cl A}}\wt{H}_{\wt{\cl B}}, \wt{\varphi}_A, 
 \wt{\varphi}_B, \wt{\xi})$ be quantum commuting models
 over $(\cl S_A,\cl S_B)$. 
 Suppose that $S\preceq_{\rm a} \wt{S}$ and that $\wt{S}$ is centrally supported. 
 Then the states 
 $\omega_\xi\circ\pi_{\varphi_A\cdot\varphi_B}$
 and $\omega_{\wt{\xi}}\circ\pi_{\wt{\varphi}_A\cdot\wt{\varphi}_B}$ of $C^*_u(\cl S_A)\otimes_{\max} C^*_u(\cl S_B)$ coincide. 
\end{lemma}

\begin{proof}
We write $\nph = \nph_A\cdot\nph_B$ and 
$\wt{\nph} = \wt{\nph}_A\cdot \wt{\nph}_B$ for brevity. 
It suffices to show that  $(\omega_\xi\circ\pi_{\varphi})(a\otimes b) = (\omega_{\wt{\xi}}\circ\pi_{\wt{\varphi}})(a\otimes b)$ holds for all 
elements $a\in C_u^*(S_A)$ (resp. $b\in C_u^*(S_B)$) of the form 
$a = s_1\ldots s_n$ (resp. $b = t_1\ldots t_n$), where $s_i\in \cl S_A$
(resp. $t_i\in\cl S_B$), $i = 1,\dots,n$. 
Write $\pi_A(a) = \pi_\varphi(a\otimes 1)$, $\wt\pi_A(a) = \pi_{\wt\varphi}(a\otimes 1)$, 
$\pi_B(b) = \pi_\varphi(1\otimes b)$ and $\wt\pi_B(b) = \pi_{\wt\varphi}(1\otimes b)$, $a\in C_u^*(\cl S_A)$, $b\in C_u^*(\cl S_B)$. 

Let $(V_i)_{i\in \bb{N}}$ be a sequence of local isometries, 
where $V_i : H\to \wt{H}\otimes H_{i,\rm aux}$, $i\in \bb{N}$,
implementing the preorder assumption $S\preceq_{\rm a} \wt{S}$
(see Definition \ref{d_appdil} (ii)), and write
$V_i = V_{i,2,2}\circ V_{i,1,2} = V_{i,2,1}\circ V_{i,1,1}$, where 
$V_{i,1,2}$, $V_{i,2,1}$ are $\cl{B}$-local (and $V_{i,1,1}$, $V_{i,2,2}$  are $\cl{A}$-local).

We have that 
\begin{equation}\label{eq_Vi22}
V_{i,1,2}\pi_A(a) =\sigma_{A,i}(\pi_A(a))V_{i,1,2}\ \mbox{ and } \ 
V_{i,2,2}\sigma_{A,i}(\pi_A(a))= \rho_{A,i}(\pi_A(a))V_{i,2,2},
\end{equation}
for $a\in C_u^*(\cl S_A)$ and 
for some normal *-representations $\sigma_{A,i}$ and $\rho_{A,i}$ of $\cl A$, the latter mapping into $\pi_{\wt{H}}(\wt{\cl A})\bar\otimes\pi_{H_{i,\rm aux}}(\cl A_{i,\rm aux})$.
Let $\wt{\epsilon}_{A}$ be the support projection of 
the restriction of $\omega_{\wt{\xi}}$ to $\pi_{\wt{H}}({\wt{\cl A}})$
and $\wt{p}_A$ be the projection onto $\overline{\pi_{\wt{H}}(\wt{\cl A})\wt{\xi}}$. 
Set $\wt{q}_A = \wt{\epsilon}_A\wt{p}_A$. As $\wt{S}$ is centrally supported, 
\begin{equation}\label{eq_cen!}
\wt{\epsilon}_A\wt{\pi}_A(a) = \wt{\pi}_A(a)\wt{\epsilon}_A, 
\ \ \ a\in C_u^*(\cl S_A), 
\end{equation}
and hence 
\begin{equation}\label{eq_wtq}
\wt{q}_A\wt{\pi}_A(a) = \wt{\pi}_A(a)\wt{q}_A, \ \ \ a\in C_u^*(\cl S_A).
\end{equation}
We claim that for every $n\in\N$, contractions $s_1,\dots,s_n\in\cl S_A$, $\epsilon>0$, there exists $i_{0}$ such that, 
setting $a=s_1\cdots s_n$, if $i\geq i_{0}$ then
\begin{eqnarray}\label{eq_aimind}
\wt{\pi}_A(a)\wt{\xi}\otimes\xi_{i,{\rm aux}}
& \sim^{\epsilon} & (\wt{\epsilon}_A
\otimes I_{H_{i,{\rm aux}}})\rho_{A,i}(\pi_A(a))(\wt{\xi}\otimes
\xi_{i,{\rm aux}}).
\end{eqnarray}

Let $n=1$, $s\in\cl S_A$ be contractive and $\epsilon > 0$. By the relation $S\preceq_{\rm a} \wt{S}$ and (\ref{eq_Vi22}), there exists $i_0$ such that, if $i\geq i_0$ then 
\begin{eqnarray*}
\wt{\pi}_A(s)\wt{\xi}\otimes\xi_{i,\rm aux}
& \sim^{\epsilon/2} &
V_{i,2,2}V_{i,1,2}\pi_A(s)\xi = V_{i,2,2}\sigma_{A,i}(\pi_A(s))V_{i,1,2}\xi\\
& = & \rho_{A,i}(\pi_A(s))V_i\xi\\
& \sim^{\epsilon/2}&
\rho_{A,i}(\pi_A(s))(\wt{\xi}\otimes\xi_{i,\rm aux}).
\end{eqnarray*} 
Claim (\ref{eq_aimind}) now follows from (\ref{eq_supre}) and (\ref{eq_cen!}).

Assume the claim holds for $n>1$. Consider $a = s_1\ldots s_ns_{n+1}$, where each $s_{i}$ is a contraction in $\cl S_A$, and let $\epsilon>0$. By hypothesis (and the case $n=1$), there is an $i_0$ such that for every $i\geq i_0$,
\begin{eqnarray}\label{eq_aimind2}
\wt{\pi}_A(a')\wt{\xi}\otimes\xi_{i,{\rm aux}}
& \sim^{\epsilon/5} & (\wt{\epsilon}_A\otimes 1_{H_{i,{\rm aux}}})\rho_{A,i}(\pi_A(a'))(\wt{\xi}\otimes
\xi_{i,{\rm aux}}),
\end{eqnarray}
where $a'=s_1\ldots s_n$ or $a'=s_{n+1}$. By Lemmas \ref{uptoeps} and \ref{sep_cycl}, 
 there exists 
 $\tilde a\in(\wt{q}_A\pi_{\wt{H}}(\wt{\cl A})\wt{q}_A)'\subseteq 
 \cl B(\wt{q}_A\wt{H})$ such that  
$$\wt{\pi}_A(s_{n+1})\wt{\xi}=\wt{q}_A\wt{\pi}_A(s_{n+1})\wt{q}_A\wt{\xi}\sim^{\epsilon/5}\tilde a\wt{\xi}.$$
%\marginpar{\tiny L. it seems we need $\tilde a\in\pi_{\wt{H}}(\wt{\cl A})'$, but does centrally supported imply its existence, satisfying things below?}  
Using (\ref{eq_Vi22}), (\ref{eq_cen!}), (\ref{eq_wtq}), 
(\ref{eq_aimind2}) and 
the fact that 
$$\rho_{A,i}(\pi_A(a'))\in\pi_{\wt{H}}(\wt{\cl A})\bar\otimes\pi_{H_{i,\rm aux}}(\cl A_{i,\rm aux}),$$ 
for $i\geq i_0$, we have
\begin{eqnarray*}
 &&
 \hspace{-0.7cm} \wt{\pi}_A(a')\wt{\pi}_A(s_{n+1})\wt{\xi}\otimes\xi_{i,\rm aux}
 \sim^{\epsilon/5}   
 \wt{\pi}_A(a')\tilde a\wt{\xi}\otimes\xi_{i,\rm aux}\\
 && =
 \wt{q}_A\wt{\pi}_A(a')\wt{q}_A\tilde a\wt{\xi}\otimes\xi_{i,\rm aux}
 =
 (\tilde a\otimes 1_{i,\rm aux})(\wt{\pi}_A(a')\wt{\xi}\otimes\xi_{i,\rm aux})\\
 && \sim^{\epsilon/5}
(\tilde a\otimes 1_{i,\rm aux})(\wt{\epsilon}_A\otimes 1_{H_{i,\rm aux}})\rho_{A,i}(\pi_A(a'))(\wt{\xi}\otimes\xi_{i,\rm aux})\\
 && =
 (\wt{q}_A\otimes 1_{H_{i,{\rm aux}}})\rho_{A,i}(\pi_A(a'))
 (\tilde a\wt{\xi}\otimes\xi_{i,{\rm aux}})\\
 && \sim^{\epsilon/5} 
 (\wt{q}_A\otimes 1_{H_{i,\rm aux}})\rho_{A,i}(\pi_A(a'))(\wt{\pi}_A(s_{n+1})\wt{\xi}\otimes\xi_{i,{\rm aux}})\\
 && \sim^{\epsilon/5}
 (\wt{q}_A\otimes 1_{H_{i,\rm aux}})\rho_{A,i}(\pi_A(a'))V_{i,2,2}V_{i,1,2}\pi_A(s_{n+1})\xi\\
 && =
 (\wt{q}_A\otimes 1_{H_{i,\rm aux}})\rho_{A,i}(\pi_A(a'))(V_{i,2,2}\sigma_{A,i}(\pi_A(s_{n+1}))V_{i,1,2}\xi\\
 && =
 (\wt{q}_A\otimes 1_{H_{i,{\rm aux}}})\rho_{A,i}(\pi_A(a's_{n+1}))V_{i,2,2}V_{i,1,2}\xi\\
 && =
 (\wt{q}_A\otimes 1_{H_{i,\rm aux}})\rho_{A,i}(\pi_A(a))V_i\xi\\
 && \sim^{\epsilon/5}
 (\wt{\epsilon}_A\otimes 1_{H_{i,{\rm aux}}})\rho_{A,i}(\pi_A(a))(\wt{\xi}\otimes\xi_{i, {\rm aux}}).
\end{eqnarray*}
Thus, 
$$\wt{\pi}_A(a)\wt{\xi}\otimes\xi_{i,{\rm aux}} 
\sim^{\epsilon} (\wt{\epsilon}_A\otimes 1_{H_{i,{\rm aux}}})
\rho_{A,i}(\pi_A(a))(\wt{\xi}\otimes\xi_{i,{\rm aux}}).$$
A similar induction argument shows that, 
if $\wt{\epsilon}_B$ is the support projection of the restriction of  
the functional $\omega_{\wt{\xi}}$ on $\pi_{\wt{H}}(\cl B^o)$, then for all monomials $b\in C_u^*(\cl S_B)$ whose individual terms are contractions in $\cl S_B$, and $\epsilon>0$, there is an $i_0$ such that for $i\geq i_0$,
\begin{equation}\label{eq_forB}
\wt{\pi}_B(b)\wt{\xi}\otimes\xi_{i,{\rm aux}}
\sim^{\epsilon} 
(\wt{\epsilon}_B\otimes 1_{H_{i,{\rm aux}}})\rho_{B,i}(\pi_B(b))(\wt{\xi}\otimes\xi_{i,{\rm aux}}),
\end{equation}
where $\rho_{B,i}$ is the normal $*$-representation of $\cl B^o$ constructed from the locality of $V_i$, analogous to $\rho_{A,i}$.
We have that 
$\wt{\epsilon}_A\in \pi_{\wt{H}}(\wt{\cl A})$, $\wt{\epsilon}_B\in \pi_{\wt{H}}(\wt{\cl B}^o)$ and $[\pi_{\wt{H}}(x), \pi_{\wt{H}}(y^o)]=0$, $x\in \wt{\cl A}$, $y\in\wt{\cl B}$.
Using (\ref{eq_supre}), (\ref{eq_cen!}),  (\ref{eq_aimind}), (\ref{eq_forB}), the relations
$$[\wt{\pi}_A(a)\otimes 1_{H_{i,\rm aux}},\rho_{B,i}(\pi_{B}(b))] = 0$$ 
and the fact that $V_{i,2,2}V_{i,1,2} = V_{i,2,1}V_{i,1,1}$, 
for monomials $a\in C_u^*(\cl S_A)$ and $b\in C_u^*(\cl S_B)$, whose individual terms are contractions in $\cl S_A$ and $\cl S_B$, respectively and $\epsilon>0$, there is an $i_0$ such that for $i\geq i_0$, we obtain 
\begin{eqnarray*}
    &&
    (\omega_{\wt{\xi}}\circ\pi_{\wt{\varphi}})(a\otimes b)
    =
    \langle \wt{\pi}_A(a)\wt{\pi}_B(b)\wt{\xi},\wt{\xi}\rangle
    =
    \langle \wt{\pi}_A(a)\wt{\pi}_B(b)\wt{\xi}\otimes\xi_{i,{\rm aux}},\wt{\xi}\otimes\xi_{i,{\rm aux}}\rangle\\
    && \sim^{\epsilon/4}
    \langle (\wt{\pi}_A(a)\wt{\epsilon}_B\otimes 1_{H_{i,{\rm aux}}})\rho_{B,i}(\pi_B(b))(\wt{\xi}\otimes\xi_{i,{\rm aux}}), \wt{\xi}\otimes\xi_{i,{\rm aux}}\rangle\\
    && =
    \langle \rho_{B,i}(\pi_B(b))(\wt{\pi}_A(a)\wt{\xi}\otimes\xi_{i,{\rm aux}}), \wt{\xi}\otimes\xi_{i,{\rm aux}}\rangle\\
    && \sim^{\epsilon/4}
    \langle \rho_{B,i}(\pi_B(b))
    (\wt{\epsilon}_A\otimes 1_{H_{i,{\rm aux}}})\rho_{A,i}(\pi_A(a))
    (\wt{\xi}\otimes\xi_{i,{\rm aux}}), 
    \wt{\xi}\otimes\xi_{i,{\rm aux}}\rangle\\
    && \sim^{\epsilon/4}
    \langle \rho_{B,i}(\pi_B(b))\rho_{A,i}(\pi_A(a))V_{i,2,2}V_{i,1,2}\xi, \wt{\xi}\otimes\xi_{i,{\rm aux}}\rangle\\
    && =
    \langle \rho_{B,i}(\pi_B(b))V_{i,2,2}V_{i,1,2}\pi_A(a)\xi, \wt{\xi}\otimes\xi_{i,{\rm aux}}\rangle\\
    && =
    \langle \rho_{B,i}(\pi_B(b))V_{i,2,1}V_{i,1,1}\pi_A(a)\xi, \wt{\xi}\otimes\xi_{i,{\rm aux}}\rangle\\
    && =
    \langle V_{i,2,1}V_{i,1,1} \pi_B(b)\pi_A(a)\xi, \wt{\xi}\otimes\xi_{i,{\rm aux}}\rangle\\
    && = 
    \langle \pi_B(b)\pi_A(a)\xi, V_i^*(\wt{\xi}\otimes\xi_{i,{\rm aux}})\rangle\\
    && \sim^{\epsilon/4}
    \langle \pi_B(b)\pi_A(a)\xi, \xi\rangle
    = 
    (\omega_\xi\circ\pi_\varphi)(a\otimes b).
    \end{eqnarray*}
Since $\epsilon$ is an arbitrary positive real, we conclude that 
%\begin{equation}\label{eq_monco}
$$(\omega_{\wt{\xi}}\circ\pi_{\wt{\varphi}})(a\otimes b) = (\omega_\xi\circ\pi_\varphi)(a\otimes b).$$
%\end{equation}
By linearity, density and continuity it follows that $\omega_\xi\circ\pi_{\varphi_A\cdot\varphi_B}
 = \omega_{\wt{\xi}}\circ\pi_{\wt{\varphi}_A\cdot\wt{\varphi}_B}$.
\end{proof}

\end{document}
